\documentclass[11pt,a4paper]{article}
\usepackage[margin=25mm]{geometry}
\usepackage[T1]{fontenc}
\usepackage{mathptmx,booktabs,longtable,hyperref}
\setlength{\parindent}{0pt}
\setlength{\parskip}{6pt}
\begin{document}
\begin{center}\Large\bfseries Supplementary Reference Paths and Evaluation Metadata\end{center}
\textbf{AD-STGN}\par
These definitions accompany the frozen v3.6 evaluation. They are exported from \texttt{tep\_meta\_final.py}; the machine-readable CSV and JSON files are supplied alongside this document. The held-out target is XMEAS(35), represented at the unit level by PURGE.
\section*{S1. Ordered reference paths}
The routes below define the reference unit sequences for ordered path recovery. They are task-specific evaluation annotations, not trajectories inferred by the model.
\begin{center}\small
\begin{tabular}{ll}\toprule Source unit & Ordered reference route \\\midrule
FEED & FEED $\rightarrow$ REACTOR $\rightarrow$ SEPARATOR $\rightarrow$ PURGE \\
RCOOL & RCOOL $\rightarrow$ REACTOR $\rightarrow$ SEPARATOR $\rightarrow$ PURGE \\
SCOOL & SCOOL $\rightarrow$ SEPARATOR $\rightarrow$ PURGE \\
REACTOR & REACTOR $\rightarrow$ SEPARATOR $\rightarrow$ PURGE \\
\bottomrule\end{tabular}\end{center}
\section*{S2. Unit-level process-influence reference graph}
These 13 directed edges are used for distinct-edge evaluation. The category labels follow the frozen metadata. The graph includes indirect process influences and must not be interpreted as a list of immediate physical piping connections. A row denotes influence from the source unit to the target unit.
\begin{center}\small
\begin{tabular}{lll}\toprule Source & Target & Category \\\midrule
FEED & REACTOR & direct material \\
REACTOR & SEPARATOR & direct material \\
SEPARATOR & COMPRESSOR & direct material \\
COMPRESSOR & REACTOR & direct material \\
SEPARATOR & PURGE & direct material \\
SEPARATOR & STRIPPER & direct material \\
SCOOL & SEPARATOR & energy coupling \\
RCOOL & REACTOR & energy coupling \\
SEPARATOR & REACTOR & recycle feedback \\
COMPRESSOR & PURGE & recycle feedback \\
REACTOR & PURGE & indirect propagation \\
FEED & COMPRESSOR & indirect propagation \\
RCOOL & STRIPPER & indirect propagation \\
\bottomrule\end{tabular}\end{center}
The separately supplied \texttt{tep\_search\_unit\_edges.csv} contains the ten allowed transitions retained in the search/support metadata. It is not interchangeable with the 13-edge evaluation graph.
\clearpage
\section*{S3. Fault scenarios and source-unit labels}
Ten held-out realizations of each scenario form the 130-event evaluation partition. Source labels are defined at the physical-unit level.
\begin{longtable}{lp{10cm}l}\toprule Fault & Programmed disturbance & Source \\\midrule\endhead
IDV(1) & A/C feed ratio, B composition constant & FEED \\
IDV(2) & B composition, A/C ratio constant & FEED \\
IDV(3) & D feed temperature & FEED \\
IDV(4) & Reactor cooling water inlet temperature & RCOOL \\
IDV(5) & Condenser cooling water inlet temperature & SCOOL \\
IDV(6) & A feed loss & FEED \\
IDV(7) & C header pressure loss & FEED \\
IDV(8) & A/B/C feed composition & FEED \\
IDV(10) & C feed temperature & FEED \\
IDV(11) & Reactor cooling water inlet temperature (random) & RCOOL \\
IDV(13) & Reaction kinetics & REACTOR \\
IDV(14) & Reactor cooling water valve sticking & RCOOL \\
IDV(15) & Condenser cooling water valve sticking & SCOOL \\
\bottomrule\end{longtable}
\section*{S4. Provenance and conventions}
Node indices in the sensor mapping are zero-based; XMEAS and XMV identifiers retain the original one-based numbering. XMEAS(35) is excluded from the 40 graph variables. Directed CSV edges use source-to-target ordering; adjacency matrices use row $i$, column $j$ for the edge $j\to i$.

Process reference: J.~J.~Downs and E.~F.~Vogel. A plant-wide industrial process control problem. \emph{Computers \& Chemical Engineering}, 17(3), 245--255 (1993). \url{https://doi.org/10.1016/0098-1354(93)80018-I}.

The unit aggregation, ordered routes, and process-influence categories are the task-specific definitions supplied with this implementation. This export does not constitute additional independent physical validation.
\end{document}
